\documentclass{article}

% if you need to pass options to natbib, use, e.g.:
%     \PassOptionsToPackage{numbers, compress}{natbib}
% before loading neurips_2026

% The authors should use one of these tracks.
% Before accepting by the NeurIPS conference, select one of the options below.
% 0. "default" for submission
\usepackage[preprint]{neurips_2026}
% the "default" option is equal to the "main" option, which is used for the Main Track with double-blind reviewing.
% 1. "main" option is used for the Main Track
%  \usepackage[main]{neurips_2026}
% 2. "position" option is used for the Position Paper Track
%  \usepackage[position]{neurips_2026}
% 3. "eandd" option is used for the Evaluations & Datasets Track
 % \usepackage[eandd]{neurips_2026}
% 4. "creativeai" option is used for the Creative AI Track
%  \usepackage[creativeai]{neurips_2026}
% 5. "sglblindworkshop" option is used for the Workshop with single-blind reviewing
 % \usepackage[sglblindworkshop]{neurips_2026}
% 6. "dblblindworkshop" option is used for the Workshop with double-blind reviewing
%  \usepackage[dblblindworkshop]{neurips_2026}

% After being accepted, the authors should add "final" behind the track to compile a camera-ready version.
% 1. Main Track
 % \usepackage[main, final]{neurips_2026}
% 2. Position Paper Track
%  \usepackage[position, final]{neurips_2026}
% 3. Evaluations & Datasets Track
 % \usepackage[eandd, final]{neurips_2026}
% 4. Creative AI Track
%  \usepackage[creativeai, final]{neurips_2026}
% 5. Workshop with single-blind reviewing
%  \usepackage[sglblindworkshop, final]{neurips_2026}
% 6. Workshop with double-blind reviewing
%  \usepackage[dblblindworkshop, final]{neurips_2026}
% Note. For the workshop paper template, both \title{} and \workshoptitle{} are required, with the former indicating the paper title shown in the title and the latter indicating the workshop title displayed in the footnote.
% For workshops (5., 6.), the authors should add the name of the workshop, "\workshoptitle" command is used to set the workshop title.
% \workshoptitle{WORKSHOP TITLE}

% "preprint" option is used for arXiv or other preprint submissions
 % \usepackage[preprint]{neurips_2026}

% to avoid loading the natbib package, add option nonatbib:
%    \usepackage[nonatbib]{neurips_2026}

\usepackage[utf8]{inputenc} % allow utf-8 input
\usepackage[T1]{fontenc}    % use 8-bit T1 fonts
\usepackage{hyperref}       % hyperlinks
\usepackage{url}            % simple URL typesetting
\usepackage{booktabs}       % professional-quality tables
\usepackage{amsfonts}       % blackboard math symbols
\usepackage{nicefrac}       % compact symbols for 1/2, etc.
\usepackage{microtype}      % microtypography
\usepackage{xcolor}         % colors
\usepackage{algorithm}
\usepackage{algpseudocode}
\usepackage{amsmath}
\usepackage{graphicx}
\usepackage{subcaption}

% Note. For the workshop paper template, both \title{} and \workshoptitle{} are required, with the former indicating the paper title shown in the title and the latter indicating the workshop title displayed in the footnote. 
\title{The Landscape of Deep RL: Four Algorithms, Five Environments}


% The \author macro works with any number of authors. There are two commands
% used to separate the names and addresses of multiple authors: \And and \AND.
%
% Using \And between authors leaves it to LaTeX to determine where to break the
% lines. Using \AND forces a line break at that point. So, if LaTeX puts 3 of 4
% authors names on the first line, and the last on the second line, try using
% \AND instead of \And before the third author name.


\author{%
  Pakawat Nutthanithipat \\
  SCIPER: 423110 \\
  \And
  Balázs Peisz \\
  SCIPER: 390226 \\
  \And
  Andreas Wendelboe \\
  SCIPER: 423025 \\
}


\begin{document}


\maketitle


\section{Introduction}

In this report we re-implement and empirically compare four widely cited algorithms: DQN \citep{mnih2013playing}, PPO \citep{schulman2017proximal}, SAC \citep{haarnoja2018soft}, and TD3 \citep{fujimoto2018addressing}. They span value-based, on-policy policy-gradient, and off-policy actor-critic methods, and covering both discrete and continuous action spaces. We evaluate them on a set of classic OpenAI Gym control tasks (CartPole, Acrobot, MountainCar, MountainCarContinuous, and Pendulum).

\section{Environments}

All five environments used in this report belong to the \emph{Classic Control} suite of the Gymnasium library \citep{towers2024gymnasium}, a maintained fork of OpenAI Gym. They cover many properties relevant to evaluating RL algorithms: discrete versus continuous action spaces, dense versus sparse rewards, low- versus higher-dimensional state spaces, and stable versus unstable dynamics.

CartPole-v1 \citep{barto1983neuronlike} requires keeping a pole upright on a cart by pushing left or right (discrete, 2 actions) given a 4-dimensional state. A reward of $+1$ per step makes the return equal the episode length, capped at the 500-step ``solved'' threshold. It is the easiest environment in our benchmark: dense rewards, a small state space, and a near-optimal linear policy.

MountainCar-v0 \citep{moore1990efficient} places a car in a valley that must reach a flag on the right hill, but its engine is too weak to climb directly; the agent must build momentum by rocking between the hills. The state is 2-dimensional (position, velocity) and the action is discrete (accelerate left, none, right), with a reward of $-1$ per step and a 200-step truncation.

MountainCarContinuous-v0 shares the physical setup of MountainCar-v0 but replaces the discrete action with a continuous force in $[-1, 1]$. Its reward also differs: a per-step penalty of $-0.1\,a_t^2$ discourages large actions, while a one-off $+100$ bonus is granted on reaching the goal, over a 999-step horizon.

Pendulum-v1 is the classical swing-up problem: starting from a random orientation, the agent applies a continuous torque in $[-2, 2]$ N·m to bring the pendulum upright and hold it there. The 3-dimensional state encodes the angle as $(\cos\theta, \sin\theta)$ together with the angular velocity, and the per-step reward quadratically penalizes deviation from upright, angular speed, and torque, over a 200-step horizon with no termination. The reward is dense everywhere; the challenge is the underactuated dynamics, as the maximum torque cannot lift the pendulum directly and the agent must pump energy through repeated swings before stabilizing.

Acrobot-v1 \citep{sutton1996generalization} is a two-link chain, fixed at the top with the middle joint actuated, that must swing its free tip above a target height from a downward-hanging start. The 6-dimensional state encodes both link angles (as sine/cosine pairs) and their angular velocities, and the action is a discrete torque ($-1$, $0$, or $+1$ N·m) at the joint, with a reward of $-1$ per step until the goal terminates the episode. It combines the underactuated swing-up of Pendulum with the higher-dimensional state and discrete actions of CartPole, making it a useful intermediate test case.

%-----------------------------------------------------------------------------------------------------------
\section{Algorithms}
\subsection{Deep Q-Learning (DQN)}
\label{sec:dqn}
Deep Q-Learning (DQN) is a value-based, off-policy algorithm for discrete action spaces which was introduced by researchers at DeepMind in 2013 \citep{mnih2013playing} and further developed two years later \citep{mnih2015human} as a variant of Q-learning \citep{watkins1992q} based on deep neural networks. Although DQN was not the first algorithm to implement neural networks as function approximators in a reinforcement learning setting, it was the first one to do so in an efficient, end-to-end manner, solving most of the shortcomings of previous approaches like Tesauro's TD-gammon \citep{tesauro1995temporal}. 

The main issues DQN addressed were sample correlation and the non-stationarity of the underlying distribution, which both arise from the fact that the data samples in a reinforcement learning setting are inherently correlated and the data distribution is subject to change while training the algorithm (as opposed to supervised learning where the data samples are assumed to be independent and coming from the same, fixed distribution). The algorithm circumvents these problems by implementing an experience replay buffer \citep{lin1992reinforcement} which allows the model to sample mini-batches of previous transitions randomly during the training procedure, thus eliminating the sample correlation and smoothing the data distribution over past experiences.

For the project we purposefully based our implementation on the original algorithm (outlined in Algorithm \ref{alg:dqn}) as implemented in \citep{mnih2013playing}, without the addition of a target network (put forward in \citep{mnih2015human}) in order to draw sharper contrasts with the other neural network based algorithms examined in this paper (namely TD3 and SAC). The only significant modification we made to the algorithm was the substitution of a convolutional neural network architecture for a multilayer perceptron one since our test environments returned states as numerical values as opposed to images like in the original Atari setting.

\begin{algorithm}
\caption{Deep Q-learning with Experience Replay (based on \citep{mnih2013playing})}
\label{alg:dqn}
\begin{algorithmic}
\State Initialize replay memory $D$ to capacity $N$
\State Initialize action-value function $Q$ with random weights $\theta$
\For{episode $= 1, M$}
    \State Initialize starting state $s_1$ from the environment
    \For{$t = 1, T$}
        \State With probability $\epsilon$ select a random action $a_t$
        \State otherwise select $a_t = \arg\max_a Q(s_t, a; \theta)$
        \State Execute action $a_t$ in emulator and observe reward $r_t$ and next state $s_{t+1}$
        \State Store transition $(s_t, a_t, r_t, s_{t+1})$ in $D$
        \State Sample random minibatch of transitions $(s_j, a_j, r_j, s_{j+1})$ from $D$
        \State Set $y_j = 
        \begin{cases} 
        r_j & \text{for terminal } s_{j+1}\\
        r_j + \gamma \max_{a'} Q(s_{j+1}, a'; \theta) & \text{for non-terminal } s_{j+1}
        \end{cases}$
        \State Perform a gradient descent step on $(y_j - Q(s_j, a_j\theta))^2$
    \EndFor
\EndFor
\end{algorithmic}
\end{algorithm}

\subsection{Twin Delayed Deep Deterministic Policy Gradient (TD3)}
\label{sec:td3}
Twin Delayed Deep Deterministic Policy Gradient is an off-policy actor-critic algorithm \citep{fujimoto2018addressing} for continuous action spaces which serves as an upgrade to the then state of the art actor-critic method Deep Deterministic Policy Gradient (DDPG) \citep{lillicrap2019continuouscontroldeepreinforcement}, building upon it while addressing its main shortcomings. The main problem TD3 addresses is overestimation bias which is a well documented property of Q-learning in discrete settings \citep{thrun2014issues} but also present in actor-critic methods, arising from the repeated maximization of noisy value estimates resulting in constant overestimation which is exacerbated by the bootstrapping nature of temporal difference learning. 

TD3 aims to reduce the inherent overestimation bias by the introduction of "clipped double Q-learning" which entails the usage of two independent critic networks which are used in the Bellman update by taking the minimum of their output, thus binding the estimate for the value of the future state when updating the target (line 9, Algorithm \ref{alg:td3}). In addition to the implementation of the twin critics TD3 reduces the function approximation error by the application of target networks \citep{mnih2015human} and the novel idea of delayed policy updates which means updating the actor less frequently than the critics resulting in lower variance value updates and thus higher quality policy updates (line 10-13, in Algorithm \ref{alg:td3}). Last but not least TD3 tackles the problem inherent in DDPG and all methods equipped with deterministic policies, namely their tendency to overfit to narrow peaks in the value estimate, which is solved by a form of regularisation called target policy smoothing. This works similarly to the learning update used in SARSA \citep{sutton1998reinforcement}, by adding a tiny bit of random noise to the action selected by the target policy before showing it to the critic (line 8, Algorithm \ref{alg:td3}). The method reinforces the idea that similar actions should have similar values, by enforcing the critic to evaluate a small area around the target action instead of fixating on "narrow spikes", thus smootning the value estimate.

In our implementation of TD3 we stuck to the algorithm outlined in \citep{fujimoto2018addressing} as shown in Algorithm \ref{alg:td3}.

\begin{algorithm}
\caption{Twin Delayed Deep Deterministic Policy Gradient \citep{fujimoto2018addressing}}
\label{alg:td3}
\begin{algorithmic}[1]
\State Initialize critic networks $Q_{\theta_1}$, $Q_{\theta_2}$, and actor network $\pi_\phi$ with random parameters $\theta_1, \theta_2, \phi$
\State Initialize target networks $\theta'_1 \gets \theta_1, \theta'_2 \gets \theta_2, \phi' \gets \phi$
\State Initialize replay buffer $\mathcal{B}$
\For{$t = 1$ \textbf{to} $T$}
    \State Select action with exploration noise $a \sim \pi_\phi(s) + \epsilon$, where $\epsilon \sim \mathcal{N}(0, \sigma)$, and observe reward $r$ and new state $s'$
    \State Store transition tuple $(s, a, r, s')$ in $\mathcal{B}$
    \State Sample mini-batch of $N$ transitions $(s, a, r, s')$ from $\mathcal{B}$
    \State $\tilde{a} \gets \pi_{\phi'}(s') + \epsilon$, where $\epsilon \sim \text{clip}(\mathcal{N}(0, \tilde{\sigma}), -c, c)$
    \State $y \gets r + \gamma \min_{i=1,2} Q_{\theta'_i}(s', \tilde{a})$
    \State Update critics $\theta_i \gets \arg\min_{\theta_i} N^{-1} \sum \left( y - Q_{\theta_i}(s, a) \right)^2$
    \If{$t \bmod d = 0$}
        \State Update $\phi$ by the deterministic policy gradient:
        \Statex \qquad $\nabla_\phi J(\phi) = N^{-1} \sum \nabla_a Q_{\theta_1}(s, a) \big|_{a=\pi_\phi(s)} \nabla_\phi \pi_\phi(s)$
        \State Update target networks:
        \Statex \qquad $\theta'_i \gets \tau \theta_i + (1 - \tau)\theta'_i$
        \Statex \qquad $\phi' \gets \tau \phi + (1 - \tau)\phi'$
    \EndIf
\EndFor
\end{algorithmic}
\end{algorithm}

\subsection{Soft Actor-Critic (SAC)}
\label{sec:sac}

Soft Actor-Critic \citep{haarnoja2018soft} is an off-policy actor-critic algorithm for continuous action spaces which, like TD3 (Section~\ref{sec:td3}), is trained from an experience replay buffer (Section~\ref{sec:dqn}) and uses clipped double-Q learning to curb the overestimation bias described there. What sets it apart from TD3 is that it does not only maximise the expected return but also the entropy of the policy, optimising the so-called maximum-entropy objective which augments the return with the expected policy entropy:
\begin{equation}
    J(\pi) \;=\; \sum_{t=0}^{T} \mathbb{E}_{(\mathbf{s}_t,\mathbf{a}_t)\sim \rho_\pi}
    \Big[\, r(\mathbf{s}_t,\mathbf{a}_t) \;+\; \alpha\,\mathcal{H}\!\left(\pi(\cdot\,|\,\mathbf{s}_t)\right) \,\Big],
    \label{eq:sac-objective}
\end{equation}
where the temperature $\alpha \ge 0$ controls the trade-off between reward and entropy and recovers the standard RL objective in the limit $\alpha \to 0$. Whereas TD3 and DDPG explore by adding a fixed amount of noise to a deterministic action, the entropy term gives SAC a built-in exploration mechanism that depends on the state, spreading probability mass over the near-optimal actions and resulting in policies that are more robust to modelling errors \citep{ziebart2010thesis,haarnoja2017energy}.

\paragraph{Function approximators and updates.}
SAC approximates maximum-entropy policy iteration \citep{haarnoja2018soft} with a stochastic policy $\pi_\phi$, a soft Q-function and a separate soft value function $V_\psi$. The policy is a Gaussian whose mean and log standard deviation are the outputs of an MLP, with a $\tanh$ squashing applied to the sampled action so that it stays inside the valid range, together with the corresponding change-of-variables correction to its log-probability. All networks are trained on minibatches $(\mathbf{s},\mathbf{a},r,\mathbf{s}')\sim\mathcal{D}$. The Q-function minimises a soft Bellman residual whose bootstrap target adds an entropy term to the next-state value and, just like in TD3, takes the minimum $\min_{i=1,2} Q_{\theta_i}$ over the two critics in order to limit overestimation, while the target value network is kept as a Polyak-averaged copy (Section~\ref{sec:td3}). The policy is updated by minimising the expected KL divergence to the exponential of the Q-function, and using the reparameterisation $\mathbf{a} = f_\phi(\epsilon;\mathbf{s})$ with $\epsilon$ sampled from a fixed noise distribution this turns into the low-variance objective
\begin{equation}
    J_\pi(\phi) \;=\; \mathbb{E}_{\mathbf{s}\sim\mathcal{D},\,\epsilon\sim\mathcal{N}}
    \!\left[\log \pi_\phi\!\left(f_\phi(\epsilon;\mathbf{s})\,|\,\mathbf{s}\right) -
        Q_\theta\!\left(\mathbf{s},\,f_\phi(\epsilon;\mathbf{s})\right)\right].
    \label{eq:sac-pi-loss}
\end{equation}
The full update (Algorithm~\ref{alg:sac}) alternates one environment step under the current stochastic policy with gradient steps on these losses and a Polyak update of the target network. Since the critic updates only need $(\mathbf{s},\mathbf{a},r,\mathbf{s}')$ tuples SAC is fully off-policy and can reuse the entire buffer, the only on-policy quantities being the freshly sampled actions used to form the value- and policy-loss targets.

\begin{algorithm}[t]
\caption{Soft Actor-Critic \citep{haarnoja2018soft}}
\label{alg:sac}
\begin{algorithmic}
\State Initialise parameters $\psi, \bar\psi, \theta_1, \theta_2, \phi$ and replay buffer $\mathcal{D} \gets \emptyset$.
\For{each iteration}
    \For{each environment step}
        \State $\mathbf{a}_t \sim \pi_\phi(\cdot\,|\,\mathbf{s}_t)$;\;
               $\mathbf{s}_{t+1} \sim p(\cdot\,|\,\mathbf{s}_t,\mathbf{a}_t)$
        \State $\mathcal{D} \gets \mathcal{D} \cup \{(\mathbf{s}_t,\mathbf{a}_t,r_t,\mathbf{s}_{t+1})\}$
    \EndFor
    \For{each gradient step}
        \State $\psi \gets \psi - \lambda_V \hat\nabla_\psi J_V(\psi)$
        \State $\theta_i \gets \theta_i - \lambda_Q \hat\nabla_{\theta_i} J_Q(\theta_i)\;\;\text{for } i \in \{1,2\}$
        \State $\phi \gets \phi - \lambda_\pi \hat\nabla_\phi J_\pi(\phi)$
        \State $\bar\psi \gets \tau \psi + (1-\tau)\bar\psi$
    \EndFor
\EndFor
\end{algorithmic}
\end{algorithm}

\subsection{Proximal Policy Optimization (PPO)}

% \subsubsection{Baseline Algorithm}

Proximal Policy Optimization \citep{schulman2017proximal} is an on-policy actor-critic algorithm that addresses a core difficulty in policy gradients: large update steps can collapse performance, yet explicit KL constraints (as in TRPO) are computationally expensive. PPO-Clip resolves this with a simple clipped surrogate objective that implicitly enforces a trust region.

Let $\pi_{\theta_k}$ denote the behaviour policy that collected rollout $\mathcal{D}_k$. PPO measures deviation from this policy via the probability ratio $r_t(\theta) = \pi_\theta(a_t \mid s_t) / \pi_{\theta_k}(a_t \mid s_t)$, which equals one when $\theta = \theta_k$. The clipped surrogate objective is
\begin{equation}
    L^{\mathrm{CLIP}}(\theta)
    \;=\; \mathbb{E}_{t \sim \mathcal{D}_k}
    \!\left[
        \min\!\left(
            r_t(\theta)\,\hat{A}_t,\;
            \mathrm{clip}\!\left(r_t(\theta),\,1-\varepsilon,\,1+\varepsilon\right)\hat{A}_t
        \right)
    \right],
    \label{eq:ppo-clip}
\end{equation}
where $\varepsilon$ (typically $0.2$) bounds how far the ratio can move. The $\min$ creates a lower bound: when $\hat{A}_t > 0$, the update is capped at $r_t = 1+\varepsilon$; when $\hat{A}_t < 0$, it is floored at $1-\varepsilon$.

Advantages $\hat{A}_t$ are estimated via GAE-$\lambda$ \citep{schulman2016gae}, which balances bias and variance through a weighted sum of TD residuals $\delta_t = r_t + \gamma V(s_{t+1})(1-d_t) - V(s_t)$:
\begin{equation}
    \hat{A}_t \;=\; \sum_{l=0}^{T-t-1} (\gamma\lambda)^l\,\delta_{t+l}.
    \label{eq:gae}
\end{equation}
Setting $\lambda = 1$ recovers Monte Carlo returns (high variance, low bias); $\lambda = 0$ reduces to one-step TD (low variance, higher bias). The return target is $\hat{R}_t = \hat{A}_t + V_{\theta_k}(s_t)$.

The actor and critic share a neural network with separate output heads. Each gradient step minimises
\begin{equation}
    \mathcal{L}(\theta)
    \;=\; -L^{\mathrm{CLIP}}(\theta)
    \;+\; c_v\,\mathbb{E}_t\!\left[\left(V_\theta(s_t) - \hat{R}_t\right)^2\right]
    \;-\; c_e\,\mathbb{E}_t\!\left[\mathcal{H}\!\left(\pi_\theta(\cdot \mid s_t)\right)\right],
    \label{eq:ppo-loss}
\end{equation}
where $c_v$ and $c_e$ weight the value loss and entropy bonus. The entropy term encourages exploration by penalising premature collapse of the action distribution. Algorithm~\ref{alg:ppo} summarises the training loop.

% \paragraph{Algorithm:} PPO alternates between collecting on-policy rollouts and performing multiple epochs of minibatch gradient descent on the same data (Algorithm~\ref{alg:ppo}). Unlike off-policy methods, the rollout buffer is discarded after each update cycle, so all data used for training was generated by the current policy. To keep the probability ratios $r_t(\theta)$ well-defined, the log-probabilities $\log\pi_{\theta_k}(a_t \mid s_t)$ are stored at collection time and held fixed during the $K$ update epochs. Advantages are normalised to zero mean and unit variance within each minibatch to stabilise training, and gradient norms are clipped to prevent excessively large update steps.

\begin{algorithm}[H]
\caption{Proximal Policy Optimization -- PPO-Clip \citep{schulman2017proximal}}
\label{alg:ppo}
\begin{algorithmic}
\State Initialise actor-critic parameters $\theta$ and empty rollout buffer $\mathcal{B}$.
\While{total environment steps not reached}
    \For{each environment step (up to horizon $T$)}
        \State Sample $a_t \sim \pi_\theta(\cdot \mid s_t)$;\; store $(s_t, a_t, r_t, d_t, \log\pi_\theta(a_t \mid s_t))$ in $\mathcal{B}$
    \EndFor
    \State Bootstrap $V_\theta(s_T)$;\; compute $\hat{A}_t$ and $\hat{R}_t$ via GAE-$\lambda$ (backward pass over $\mathcal{B}$)
    \For{$K$ epochs over $\mathcal{B}$}
        \For{each minibatch of size $M$}
            \State Normalise $\hat{A}_t$ to zero mean and unit variance on the minibatch
            \State Compute $r_t(\theta) = \pi_\theta(a_t \mid s_t)\,/\,\pi_{\theta_k}(a_t \mid s_t)$
            \State Evaluate $L^{\mathrm{CLIP}}$, value loss, and entropy bonus (Eq.~\ref{eq:ppo-loss})
            \State $\theta \gets \theta - \alpha\,\hat\nabla_\theta\,\mathcal{L}(\theta)$;\; clip $\|\nabla_\theta \mathcal{L}\|_2$ at max norm
        \EndFor
    \EndFor
    \State Discard $\mathcal{B}$ (on-policy: data is no longer valid after $\theta$ has changed)
\EndWhile
\end{algorithmic}
\end{algorithm}

% \subsubsection{Hyperparameter Optimization}
% \label{sec:ppo-hpo}

% We study PPO's sensitivity using three approaches: (i) one-factor-at-a-time (OFAT) sweeps to isolate individual effects, (ii) random search as a model-free baseline, and (iii) Optuna TPE~\citep{akiba2019optuna} as a Bayesian alternative. All methods use the same protocol: after training, we evaluate the frozen policy over 10 deterministic episodes.

% Table~\ref{tab:ppo-search-space} defines the search space. We use a conservative space for easy environments (CartPole, Acrobot) and an extended space including $\gamma$, $n_{\mathrm{steps}}$, and network size for harder tasks.

% \begin{table}[H]
%   \centering
%   \caption{PPO search space. Easy environments exclude $\gamma$, $n_{\mathrm{steps}}$, and hidden sizes.}
%   \label{tab:ppo-search-space}
%   \small
%   \begin{tabular}{lll}
%     \toprule
%     Parameter & Type & Range / Choices \\
%     \midrule
%     $\mathrm{lr}$ & log-uniform & $[5{\times}10^{-5},\, 10^{-2}]$ \\
%     $\varepsilon_{\mathrm{clip}}$ & categorical & $\{0.1, 0.2, 0.3, 0.4\}$ \\
%     $c_e$ (entropy) & log-uniform & $[10^{-4},\, 0.1]$ \\
%     $\lambda_{\mathrm{GAE}}$ & uniform & $[0.8,\, 1.0]$ \\
%     $\gamma$ & categorical & $\{0.9, 0.95, 0.99, 0.995, 0.999\}$ \\
%     $n_{\mathrm{epochs}}$ & categorical & $\{3, 5, 10, 15, 20, 30\}$ \\
%     Batch size & categorical & $\{32, 64, 128, 256, 512\}$ \\
%     $n_{\mathrm{steps}}$ & categorical & $\{512, 1024, 2048, 4096\}$ \\
%     Hidden sizes & categorical & $\{(32,32), (64,64), (128,128), (256,256)\}$ \\
%     \bottomrule
%   \end{tabular}
% \end{table}

% OFAT varies one hyperparameter while holding others at the paper default~\citep{schulman2017proximal}, averaged over 3 seeds. Random search draws $N$ independent configurations from the joint space~\citep{bergstra2012random}. Optuna TPE models the objective and proposes trials expected to improve over past observations~\citep{bergstra2011algorithms}. Algorithm~\ref{alg:hpo} summarises the search loop.

% \begin{algorithm}[H]
% \caption{Hyperparameter search (Random / Optuna)}
% \label{alg:hpo}
% \begin{algorithmic}[1]
% \Require Search space $\mathcal{H}$, number of trials $N$, training budget $T$, eval episodes $E$
% \State $\mathcal{R} \gets \emptyset$ \Comment{Results}
% \For{$i = 1, \ldots, N$}
%     \State $\mathbf{h} \gets \textsc{Sample}(\mathcal{H}, \mathcal{R})$ \Comment{Random: uniform; TPE: surrogate model}
%     \State $\pi \gets \textsc{TrainPPO}(\mathbf{h}, T)$ \Comment{Train for $T$ env steps}
%     \State $s \gets \frac{1}{E}\sum_{j=1}^{E} G_j(\pi)$ \Comment{Mean return over $E$ deterministic episodes}
%     \State $\mathcal{R} \gets \mathcal{R} \cup \{(\mathbf{h}, s)\}$
% \EndFor
% \State \Return $\arg\max_{(\mathbf{h},s)\in\mathcal{R}} s$
% \end{algorithmic}
% \end{algorithm}

\section{Results}


\subsection{Untuned Comparison}
\label{sec:vanilla-comparison}

Before tuning, we compared the algorithms on every compatible environment using each paper's default hyperparameters; the training curves are collected in Appendix~\ref{sec:vanilla-curves} (Figures~\ref{fig:compare-vanilla-cartpole}--\ref{fig:compare-vanilla-mountaincarcontinuous}). Three patterns stand out. 

First, on the easy dense-reward tasks the methods largely agree: both discrete algorithms solve Acrobot-v1 at comparable returns (Figure~\ref{fig:compare-vanilla-acrobot}, DQN converging slightly faster, PPO marginally higher at the end), and on the continuous side SAC and TD3 both solve Pendulum-v1, with SAC converging fastest (Figure~\ref{fig:compare-vanilla-pendulum}). Second, the off-policy value-based methods are markedly less stable: on CartPole-v1, PPO reaches and holds the 500-step ceiling whereas our target-network-free DQN peaks early and then degrades (Figure~\ref{fig:compare-vanilla-cartpole}), the instability anticipated in Section~\ref{sec:dqn}. Third, the sparse-reward tasks cleanly separate the algorithms by their exploration mechanism: vanilla PPO never escapes the $-200$ plateau on MountainCar-v0 while DQN eventually does (Figure~\ref{fig:compare-vanilla-mountaincar}) due to its more aggressive epsilon-greedy policy. 

In the continuous case (MountainCarContinuous-v0)(Figure~\ref{fig:compare-vanilla-mountaincarcontinuous}) none of the models manage to solve the environment with TD3 and PPO reaching a local optimum around 0 while SAC stalls at the $\approx-30$ plateau analysed in Section~\ref{sec:sac}. This is mostly due to the fact that while TD3 and PPO quickly learn that (without seeing the target early enough) the best strategy is minimizing the control actions i.e. staying put at the bottom, thus converging to the local minima, SAC is only able to achieve a slightly lower result since it does not only optimises its reward but also its entropy which encourages it more to take "unnecessary" actions.

\subsection{Tuned Comparison}

\begin{table}[H]
  \caption{Final-performance comparison across algorithms and environments. Each entry reports the mean undiscounted episodic return $\pm$ one standard deviation over 3 random seeds, evaluated on 200 episodes of the final policy. Dashes (---) indicate that the algorithm is incompatible with the environment's action space (DQN: discrete only; SAC/TD3: continuous only). Best score per environment in \textbf{bold}.}
  \label{tab:results}
  \centering
  \small
  \begin{tabular}{lcccc}
    \toprule
    Environment & DQN & PPO & SAC & TD3 \\
    \midrule
    CartPole-v1               & TODO         & $\mathbf{500.0 \pm 0.0}$  & ---                       & ---               \\
    Acrobot-v1                & $-86.4 \pm 8.2$          & $\mathbf-85.9 \pm 6.1$  & ---                       & ---               \\
    MountainCar-v0            & TODO & $-200.0 \pm 0.0$        & ---                       & ---               \\
    MountainCarContinuous-v0  & ---                      & $0.0 \pm 0.0$            & $\mathbf{95.53 \pm 1.24}$  & TODO    \\
    Pendulum-v1               & ---                      & $-975.4 \pm 73.4$         & $\mathbf{-114.55 \pm 70.38}$ & $-148.7 \pm 11.3$ \\
    \bottomrule
  \end{tabular}
\end{table}

\section{Discussion}


\subsection{Soft Actor-Critic (SAC)}

The two continuous-control environments stress SAC in opposite ways. Pendulum-v1 is essentially solved out of the box and is largely insensitive to tuning, whereas vanilla SAC fails outright on MountainCarContinuous-v0 and only starts to learn after two targeted changes.

\paragraph{MountainCarContinuous needs a warmup and discount change.}
With the default hyperparameters from \citet{haarnoja2018soft} MountainCarContinuous fails outright, across all three seeds the return stays flat at roughly $-30$ and the car never reaches the goal. Two changes fix this, both of which leave SAC's learning update (Eqs.~\ref{eq:sac-objective}--\ref{eq:sac-pi-loss} and Algorithm~\ref{alg:sac}) untouched and only affect the warmup phase that seeds the replay buffer with random actions before learning begins. The first one is to repeat each random warmup action for ten consecutive steps (\texttt{warmup\_action\_repeat}$=10$), which lets the car build up enough momentum to escape the valley so that the buffer actually contains goal-reaching trajectories. The second one is to raise the discount to $\gamma = 0.9999$, which keeps the one-off $+100$ terminal bonus from being washed out over the long path to the goal. The sensitivity sweep below confirms that both of them are essential.

\paragraph{Hyperparameter sensitivity.}
We ran 143-trial Optuna sweeps on both environments, reporting the mean undiscounted return of the final policy over three seeds. On MountainCarContinuous-v0 only about 21\% of the trials solve the task, and success hinges on exactly the two knobs above. \texttt{warmup\_action\_repeat} has to sit in a narrow band, the values 5 and 10 work while 1 never solves the task and 20 manages it only once, and $\gamma$ has to be large, almost all of the successful trials used $\gamma=0.999$ or $0.9999$ and the smaller values essentially never solve it (with only the occasional success at $\gamma=0.99$). Once the buffer holds successful trajectories SAC is largely insensitive to $\alpha$, the learning rate and $\tau$. Pendulum-v1 by contrast is much more forgiving, tuning barely moves the already near-optimal default and the swept configuration is almost indistinguishable from it across all three seeds (Figure~\ref{fig:sac-pendulum-tuning}).


\subsection{Proximal policy optimization (PPO)}

\paragraph{One-factor-at-a-time (OFAT):}
OFAT sweeps reveal that PPO's sensitivity is environment-dependent.
On \textbf{CartPole-v1} (Figure~\ref{fig:cartpole_ppo_hyperparameters}, Appendix), learning rate has minimal impact once above $10^{-4}$, and all tested $\varepsilon_{\mathrm{clip}}$ values converge to near-optimal return; the task is simple enough that most reasonable settings succeed.
On \textbf{Acrobot-v1}, moderate learning rates ($3{\times}10^{-4}$) outperform both lower and higher values, and non-zero entropy ($c_e \approx 0.01$) accelerates early exploration without harming final performance.
\textbf{MountainCar-v0} is the most sensitive: only $c_e \geq 0.01$ escapes the local minimum at $-200$; $\lambda_{\mathrm{GAE}} = 1.0$ (Monte Carlo returns) outperforms lower values; and $\gamma \geq 0.99$ is required for adequate credit assignment over the long horizon.
On \textbf{Pendulum-v1}, higher learning rates ($\geq 10^{-3}$) learn faster, and $\lambda_{\mathrm{GAE}} \in [0.9, 0.95]$ balances bias and variance effectively.

\paragraph{Joint search: random vs.\ Optuna:}
Table~\ref{tab:ppo-joint-search} summarises joint-search results on Pendulum-v1 and MountainCar-v0 (10 trials each, seed 0, $10^5$ steps).
On \textbf{Pendulum}, random search found a configuration scoring $-144.6$, substantially better than Optuna's best ($-954.8$). This counter-intuitive result arises because (i) the trial budget is small relative to the 9-dimensional space, limiting TPE's ability to build an accurate surrogate, and (ii) random search uniformly covers the space, occasionally sampling high-quality regions that TPE's early model underweights.
On \textbf{MountainCar}, only 1 of 10 random trials escaped the $-200$ failure mode, achieving $-140.4$; the remaining 9 trials failed to reach the goal within the training budget. This highlights MountainCar's sparse-reward challenge: without sufficient entropy or the right $\gamma$, PPO converges to a degenerate policy.

\begin{table}[htbp]
  \centering
  \caption{PPO joint-search results (10 trials, seed 0). Eval return = mean over 10 deterministic episodes after training. Higher is better. MountainCarContinuous results pending.}
  \label{tab:ppo-joint-search}
  \small
  \begin{tabular}{lcccc}
    \toprule
    Environment & Method & Best Eval & Median & Success Rate \\
    \midrule
    \multirow{2}{*}{Pendulum-v1} 
      & Random & $\mathbf{-144.6}$ & $-1092.6$ & -- \\
      & Optuna TPE & $-954.8$ & $-1284.5$ & -- \\
    \midrule
    MountainCar-v0 & Random & $-140.4$ & $-200.0$ & 10\% \\
    & Optuna TPE & ... & ... & ... \\
    \midrule
    MountainCarCont. & Random & \multicolumn{3}{c}{\emph{(pending)}} \\
                     & Optuna TPE & \multicolumn{3}{c}{\emph{(pending)}} \\
    \bottomrule
  \end{tabular}
\end{table}

\paragraph{Optuna importance analysis:}
Figure~\ref{fig:importances} shows fANOVA-based hyperparameter importances from the Optuna study on Pendulum-v1. Network capacity (\texttt{hidden\_sizes}), training intensity ($n_{\mathrm{epochs}}$, batch size), and temporal parameters ($\gamma$, $\mathrm{lr}$, $\lambda_{\mathrm{GAE}}$) each contribute 13--19\% of outcome variance. In contrast, $\varepsilon_{\mathrm{clip}}$ and $c_e$ contribute only 2\% each, consistent with PPO's design goal of robustness to the clipping threshold once it lies in a sensible range.

\paragraph{Learning dynamics:}
Figure~\ref{fig:best-random-vs-optuna} compares training episode returns for the best random and Optuna trials on Pendulum. The random-search winner (trial 5) shows steady improvement throughout training, while Optuna's best (trial 9) plateaus earlier at a worse level. Figures~\ref{fig:random-eval} and~\ref{fig:optuna-top3} (Appendix) plot the \emph{evaluation} metric used for scoring: 10 deterministic episodes at each checkpoint. The final dashed line ($-144.6$ vs.\ $-954.8$) is the trial score reported in Table~\ref{tab:ppo-joint-search}.

\paragraph{Best configurations:}
Table~\ref{tab:pendulum-best-hparams} lists the winning hyperparameters for Pendulum. Both methods selected $\varepsilon_{\mathrm{clip}}=0.1$, $n_{\mathrm{epochs}}=5$, $n_{\mathrm{steps}}=4096$, and $(128,128)$ hidden layers. The key difference is $\gamma$: random's best uses $\gamma=0.9$, while Optuna chose $\gamma=0.995$. Pendulum's 200-step episodes with dense rewards benefit from lower discount factors that reduce variance in the value bootstrap---a finding consistent with our SAC analysis on the same environment.

\paragraph{Limitations:}
Our HPO experiments use a single training seed and 10 evaluation episodes per trial, introducing non-negligible variance in trial rankings. Large performance gaps (e.g., $-144$ vs.\ $-955$) are unlikely to reverse, but close trials may swap order under resampling. Future work should increase seed count and trial budget for more robust conclusions.


\section{Conclusion}

\bibliographystyle{plainnat}
\bibliography{references}


%%%%%%%%%%%%%%%%%%%%%%%%%%%%%%%%%%%%%%%%%%%%%%%%%%%%%%%%%%%%

\appendix

\section{Appendix}

\subsection{Comparing Training Curves}
\subsubsection{Untuned (Vanilla) Algorithms}
\label{sec:vanilla-curves}

\begin{figure}[H]
    \centering
    \includegraphics[width=1\linewidth]{compare_vanilla_CartPole-v1.png}
    \caption{Vanilla training curves on CartPole-v1, comparing the algorithms that support a discrete action space (DQN and PPO). Each curve is the across-seed mean of the rolling-mean episode return.}
    \label{fig:compare-vanilla-cartpole}
\end{figure}

\begin{figure}[H]
    \centering
    \includegraphics[width=1\linewidth]{compare_vanilla_Acrobot-v1.png}
    \caption{Vanilla training curves on Acrobot-v1, comparing the discrete-action algorithms (DQN and PPO). Each curve is the across-seed mean of the rolling-mean episode return.}
    \label{fig:compare-vanilla-acrobot}
\end{figure}

\begin{figure}[H]
    \centering
    \includegraphics[width=1\linewidth]{compare_vanilla_MountainCar-v0.png}
    \caption{Vanilla training curves on MountainCar-v0, comparing the discrete-action algorithms (DQN and PPO). Each curve is the across-seed mean of the rolling-mean episode return.}
    \label{fig:compare-vanilla-mountaincar}
\end{figure}

\begin{figure}[H]
    \centering
    \includegraphics[width=1\linewidth]{compare_vanilla_Pendulum-v1.png}
    \caption{Vanilla training curves on Pendulum-v1, comparing the continuous-action algorithms (SAC, TD3 and PPO). Each curve is the across-seed mean of the rolling-mean episode return.}
    \label{fig:compare-vanilla-pendulum}
\end{figure}

\begin{figure}[H]
    \centering
    \includegraphics[width=1\linewidth]{compare_vanilla_MountainCarContinuous-v0.png}
    \caption{Vanilla training curves on MountainCarContinuous-v0, comparing the continuous-action algorithms (SAC, TD3 and PPO). Each curve is the across-seed mean of the rolling-mean episode return.}
    \label{fig:compare-vanilla-mountaincarcontinuous}
\end{figure}

\subsection{PPO Hyperparameter Experiments}
\subsubsection{One factor at a time}
\begin{figure}[H]
    \centering
    \begin{subfigure}[b]{0.49\textwidth}
        \centering
        \includegraphics[width=\linewidth]{Figure/PPO/CartPole-v1_1.png}
        \caption{Learning rate}
        \label{fig:cartpole_lr}
    \end{subfigure}
    \hfill
    \begin{subfigure}[b]{0.49\textwidth}
        \centering
        \includegraphics[width=\linewidth]{Figure/PPO/CartPole-v1_2.png}
        \caption{Clipping parameter ($\epsilon$)}
        \label{fig:cartpole_clip_eps}
    \end{subfigure}

    \caption{Effects of learning rate and clipping parameter on PPO performance in CartPole-v1. Results are averaged over three random seeds and smoothed using a moving average window of 10 episodes. Higher learning rates generally accelerate initial learning but may introduce greater variance, while different clipping values affect the trade-off between update stability and learning speed.}
    \label{fig:cartpole_ppo_hyperparameters}
\end{figure}

\begin{figure}[H]
    \centering
    \begin{subfigure}[b]{0.49\textwidth}
        \centering
        \includegraphics[width=\linewidth]{Figure/PPO/Acrobot-v1_1.png}
        \caption{Learning rate}
        \label{fig:acrobot_lr}
    \end{subfigure}
    \hfill
    \begin{subfigure}[b]{0.49\textwidth}
        \centering
        \includegraphics[width=\linewidth]{Figure/PPO/Acrobot-v1_2.png}
        \caption{Entropy coefficient}
        \label{fig:acrobot_ent_coef}
    \end{subfigure}

    \caption{Effects of learning rate and entropy coefficient on PPO performance in Acrobot-v1. Results are averaged over three random seeds and smoothed using a moving average window of 10 episodes. Different learning rates affect convergence speed and stability, while varying entropy coefficients influence the trade-off between exploration and exploitation during training.}
    \label{fig:acrobot_ppo_hyperparameters}
\end{figure}

\begin{figure}[H]
    \centering
    \begin{subfigure}[b]{0.49\textwidth}
        \centering
        \includegraphics[width=\linewidth]{Figure/PPO/Pendulum-v1_1.png}
        \caption{Learning rate}
        \label{fig:pendulum_lr}
    \end{subfigure}
    \hfill
    \begin{subfigure}[b]{0.49\textwidth}
        \centering
        \includegraphics[width=\linewidth]{Figure/PPO/Pendulum-v1_2.png}
        \caption{Entropy coefficient}
        \label{fig:pendulum_gae_lambda}
    \end{subfigure}

    \caption{Effects of learning rate and GAE parameter ($\lambda$) on PPO performance in Pendulum-v1. Results are averaged over three random seeds and smoothed using a moving average window of 10 episodes. The learning rate controls the magnitude of policy updates, while the GAE parameter ($\lambda$) affects the bias–variance trade-off in advantage estimation.}
    \label{fig:pendulum_ppo_hyperparameters}
\end{figure}
\begin{figure}[H]
    \centering
    \begin{subfigure}[b]{0.49\textwidth}
        \centering
        \includegraphics[width=\linewidth]{Figure/PPO/MountainCar-v0_1.png}
        \caption{Entropy coefficient}
        \label{fig:mountaincar_ent_coef}
    \end{subfigure}
    \hfill
    \begin{subfigure}[b]{0.49\textwidth}
        \centering
        \includegraphics[width=\linewidth]{Figure/PPO/MountainCar-v0_2.png}
        \caption{GAE parameter ($\lambda$)}
        \label{fig:mountaincar_gae_lambda}
    \end{subfigure}

    \vspace{0.5em}

    \begin{subfigure}[b]{0.6\textwidth}
        \centering
        \includegraphics[width=\linewidth]{Figure/PPO/MountainCar-v0_3.png}
        \caption{Discount factor ($\gamma$)}
        \label{fig:mountaincar_gamma}
    \end{subfigure}

    \caption{Effects of entropy coefficient, GAE parameter ($\lambda$), and discount factor ($\gamma$) on PPO performance in MountainCar-v0. Results are averaged over three random seeds and smoothed using a moving average window of 10 episodes. The entropy coefficient controls exploration, the GAE parameter ($\lambda$) adjusts the bias–variance trade-off in advantage estimation, and the discount factor ($\gamma$) determines the weighting of future rewards.}
    \label{fig:mountaincar_ppo_hyperparameters}
\end{figure}
\subsubsection{Random Search}

\begin{figure}[H]
  \centering
  \includegraphics[width=0.7\linewidth]{Figure/PPO/Pendulum-v1_5.png}
  \caption{Top-3 trials by final eval score.}
  \label{fig:hpo-curves_random_search}
\end{figure}

\begin{figure}[H]
  \centering
  \includegraphics[width=0.85\linewidth]{Figure/PPO/Pendulum-v1_7.png}
  \caption{Best random-search trial on \texttt{Pendulum-v1} (trial~5, seed~0).
    Evaluation return vs.\ training steps: at each checkpoint, the policy is
    frozen and tested for 10 deterministic episodes; points show the mean return.
    The dashed line is the final trial score ($-144.6$) used for HPO ranking.}
  \label{fig:random-eval}
\end{figure}


\begin{figure}[H]
  \centering
  \includegraphics[width=0.7\linewidth]{Figure/PPO/MountainCar-v0_5.png}
  \caption{Top-3 trials by final eval score.}
  \label{fig:hpo-curves_random_search}
\end{figure}
\begin{figure}[H]
  \centering
  \includegraphics[width=0.85\linewidth]{Figure/PPO/MountainCar-v0_7.png}
  \caption{Best Optuna TPE trial on \texttt{Pendulum-v1} (trial~9, seed~0).
    Same evaluation protocol as Figure~\ref{fig:random-eval}.
    Final trial score: $-954.8$.}
  \label{fig:optuna-top3}
\end{figure}

\begin{figure}[H]
  \centering
  \includegraphics[width=0.7\linewidth]{Figure/PPO/MountainCarContinuous-v0_2.png}
  \caption{Top-3 trials by final eval score.}
  \label{fig:hpo-curves_random_search}
\end{figure}
\begin{figure}[H]
  \centering
  \includegraphics[width=0.85\linewidth]{Figure/PPO/MountainCarContinuous-v0_6.png}
  \caption{Best Optuna TPE trial on \texttt{Pendulum-v1} (trial~9, seed~0).
    Same evaluation protocol as Figure~\ref{fig:random-eval}.
    Final trial score: $-954.8$.}
  \label{fig:optuna-top3}
\end{figure}


% \begin{figure}[H]
%   \centering
%   \begin{subfigure}[b]{0.49\linewidth}
%     \includegraphics[width=\linewidth]{Figure/PPO/Pendulum-v1_5.png}
%     \caption{Random search.}
%   \end{subfigure}
%   \hfill
%   \begin{subfigure}[b]{0.49\linewidth}
%     \includegraphics[width=\linewidth]{Figure/PPO/Pendulum-v1_6.png}
%     \caption{Optuna TPE.}
%   \end{subfigure}
%   \caption{Top-3 trials by final eval score.}
%   \label{fig:hpo-curves}
% \end{figure}

\subsubsection{Bayesian Search}

\begin{figure}[H]
  \centering
  \includegraphics[width=0.65\linewidth]{Figure/PPO/Pendulum-v1_3.png}
  \caption{Optuna fANOVA hyperparameter importances on Pendulum-v1.}
  \label{fig:importances}
\end{figure}

\begin{figure}[H]
  \centering
  \includegraphics[width=0.65\linewidth]{Figure/PPO/MountainCar-v0_4.png}
  \caption{Optuna fANOVA hyperparameter importances on Pendulum-v1.}
  \label{fig:importances}
\end{figure}

\begin{figure}[H]
  \centering
  \includegraphics[width=0.65\linewidth]{Figure/PPO/MountainCarContinuous-v0_1.png}
  \caption{Optuna fANOVA hyperparameter importances on Pendulum-v1.}
  \label{fig:importances}
\end{figure}

\begin{figure}[H]
  \centering
  \includegraphics[width=0.7\linewidth]{Figure/PPO/Pendulum-v1_6.png}
  \caption{Top-3 trials by final eval score.}
  \label{fig:hpo-curves_tpe}
\end{figure}

\begin{figure}[htbp]
  \centering
  \includegraphics[width=0.85\linewidth]{Figure/PPO/Pendulum-v1_8.png}
  \caption{Best Optuna TPE trial on \texttt{Pendulum-v1} (trial~9, seed~0).
    Same evaluation protocol as Figure~\ref{fig:random-eval}.
    Final trial score: $-954.8$.}
  \label{fig:optuna-top3}
\end{figure}

\begin{figure}[H]
  \centering
  \includegraphics[width=0.7\linewidth]{Figure/PPO/MountainCarContinuous-v0_3.png}
  \caption{Top-3 trials by final eval score.}
  \label{fig:hpo-curves_tpe}
\end{figure}

\begin{figure}[htbp]
  \centering
  \includegraphics[width=0.85\linewidth]{Figure/PPO/MountainCarContinuous-v0_5.png}
  \caption{Best Optuna TPE trial on \texttt{Pendulum-v1} (trial~9, seed~0).
    Same evaluation protocol as Figure~\ref{fig:random-eval}.
    Final trial score: $-954.8$.}
  \label{fig:optuna-top3}
\end{figure}


\begin{figure}[H]
  \centering
  \includegraphics[width=0.7\linewidth]{Figure/PPO/MountainCar-v0_6.png}
  \caption{Top-3 trials by final eval score.}
  \label{fig:hpo-curves_tpe}
\end{figure}

\begin{figure}[htbp]
  \centering
  \includegraphics[width=0.85\linewidth]{Figure/PPO/MountainCar-v0_8.png}
  \caption{Best Optuna TPE trial on \texttt{Pendulum-v1} (trial~9, seed~0).
    Same evaluation protocol as Figure~\ref{fig:random-eval}.
    Final trial score: $-954.8$.}
  \label{fig:optuna-top3}
\end{figure}


\begin{figure}[H]
  \centering
  \includegraphics[width=0.7\linewidth]{Figure/PPO/MountainCarContinuous-v0_3.png}
  \caption{Top-3 trials by final eval score.}
  \label{fig:hpo-curves_tpe}
\end{figure}

\begin{figure}[htbp]
  \centering
  \includegraphics[width=0.85\linewidth]{Figure/PPO/MountainCarContinuous-v0_5.png}
  \caption{Best Optuna TPE trial on \texttt{Pendulum-v1} (trial~9, seed~0).
    Same evaluation protocol as Figure~\ref{fig:random-eval}.
    Final trial score: $-954.8$.}
  \label{fig:optuna-top3}
\end{figure}



\begin{figure}[H]
  \centering
  \includegraphics[width=0.95\linewidth]{Figure/PPO/Pendulum-v1_4.png}
  \caption{Learning curves for the best joint-search trial under random search
    and Optuna TPE on \texttt{Pendulum-v1} (seed~0).
    Faint lines: per-episode returns; bold lines: smoothed mean
    (window~10). Final selection uses mean eval return over the training run.}
  \label{fig:best-random-vs-optuna}
\end{figure}

\begin{figure}[H]
  \centering
  \includegraphics[width=0.95\linewidth]{Figure/PPO/MountainCar-v0_9.png}
  \caption{Learning curves for the best joint-search trial under random search
    and Optuna TPE on \texttt{Pendulum-v1} (seed~0).
    Faint lines: per-episode returns; bold lines: smoothed mean
    (window~10). Final selection uses mean eval return over the training run.}
  \label{fig:best-random-vs-optuna}
\end{figure}

\begin{figure}[H]
  \centering
  \includegraphics[width=0.95\linewidth]{Figure/PPO/MountainCarContinuous-v0_7.png}
  \caption{Learning curves for the best joint-search trial under random search
    and Optuna TPE on \texttt{Pendulum-v1} (seed~0).
    Faint lines: per-episode returns; bold lines: smoothed mean
    (window~10). Final selection uses mean eval return over the training run.}
  \label{fig:best-random-vs-optuna}
\end{figure}

\begin{table}[H]
  \centering
  \caption{Best search configuration per method on \texttt{Pendulum-v1}
    (10 trials, seed~0, $10^5$ env steps). Higher eval return is better.}
  \label{tab:pendulum-best-hparams}
  \begin{tabular}{lrr}
    \toprule
    Hyperparameter & Random (trial 5) & Optuna TPE (trial 9) \\
    \midrule
    Eval return      & $-144.6$  & $-954.8$ \\
    $ \mathrm{lr} $  & $2.09\times 10^{-3}$ & $3.39\times 10^{-3}$ \\
    $ \mathrm{clip\_eps} $ & 0.1 & 0.1 \\
    $ \mathrm{ent\_coef} $ & $9.63\times 10^{-3}$ & $5.43\times 10^{-3}$ \\
    $ \mathrm{gae\_lambda} $ & 0.955 & 0.807 \\
    $ \gamma $       & 0.9   & 0.995 \\
    $ n_{\mathrm{epochs}} $ & 5 & 5 \\
    Batch size       & 32    & 64 \\
    $ n_{\mathrm{steps}} $   & 4096  & 4096 \\
    Hidden sizes     & $(128,128)$ & $(128,128)$ \\
    \bottomrule
  \end{tabular}
\end{table}

\subsubsection{Best Tuning Sensitivity}

\begin{figure}[H]
  \centering
  \includegraphics[width=1\linewidth]{Figure/PPO/Pendulum-v1_9.png}
  \caption{Optuna fANOVA hyperparameter importances on Pendulum-v1.}
  \label{fig:importances}
\end{figure}

\begin{figure}[H]
  \centering
  \includegraphics[width=1\linewidth]{Figure/PPO/MountainCar-v0_10.png}
  \caption{Optuna fANOVA hyperparameter importances on Pendulum-v1.}
  \label{fig:importances}
\end{figure}

\subsection{SAC Hyperparameter Experiments}

\begin{figure}[H]
    \centering
    \includegraphics[width=1\linewidth]{pendulum_vanilla_vs_tuned.png}
    \caption{SAC on Pendulum-v1 with default versus tuned hyperparameters: mean training return over three seeds (shaded region $\pm 1$ std). The two curves are nearly indistinguishable, indicating that hyperparameter tuning yields little improvement on this dense-reward task.}
    \label{fig:sac-pendulum-tuning}
\end{figure}


%%%%%%%%%%%%%%%%%%%%%%%%%%%%%%%%%%%%%%%%%%%%%%%%%%%%%%%%%%%%

\newpage


\end{document}